\input{estilo}
\begin{document}
\begin{figura}\begin{tikzpicture}
  % la figura 2.1 del Saltzer, tal cual
  \node[fill=relleno, align=left, font=\small, inner sep=8pt, anchor=north west] at (0,0) {%
    Hardware memory devices:\\
    \quad RAM chip\\ \quad Flash memory\\ \quad Magnetic tape\\
    \quad Magnetic disk\\ \quad CD-R and DVD-R\\[2pt]
    Higher level memory systems:\\
    \quad RAID\\ \quad File system\\ \quad Database management system};
  % la interfaz que anotó el profesor
  \node[font=\ttfamily\small, anchor=north west, align=left] at (6.8,-0.3) {%
    write(nombre, valor)\\[8pt]
    valor \tikz[baseline=-0.55ex]\draw[<-] (0,0) -- (0.5,0); read(nombre)};
\end{tikzpicture}\end{figura}
\end{document}
